%\documentclass[10pt,respuestas,a4]{examenEPN}
\documentclass[10pt,a4]{examenEPN0}
\examen{Examen 01}
\usepackage{comandosEPN}

\usepackage{empheq}

\newcommand{\itemcp}{\stepcounter{enumi}\item[ \textcolor{red}{\theenumi.$^*$}] }

\newcommand{\iitemcp}{\stepcounter{enumii}\item[ \textcolor{red}{\theenumii)$^*$}] }


\materia{Álgebra Lineal}
%\autor{Departamento de Formación Básica}
\autor{GR \hspace{1cm}.}
%\fecha{Semestre 2026-A}
\fecha{Nombre: \hspace{9cm} Cédula:}
\logouno[1.6cm]{Logo_EPN.eps}
\logodos[2.6cm]{Logo_DFB.eps}

\newcommand{\la}{\lambda}

%Datos
\newcommand{\datosAlgebra}[2]{
\EnExamen
{
\vspace{5mm}
    \noindent
    Nombres: \hrulefill\
    Carrera: \rule{4.5cm}{0.4pt}\
    \ifthenelse{\equal{#1}{}}%
    {Paralelo: \underline{\parbox{1.8cm}{\parbox{1cm}{}}}\\[3mm]}
    {Paralelo: \underline{\parbox{1.8cm}{\centering #1}}\\[3mm]}
    C\'edula: \rule{5cm}{0.4pt}\
    Firma: \hrulefill\
    Fecha: \underline{\parbox{2.1cm}{\centering #2}}
}}


% \newcommand{\indicacionesPrueba}{%
% \EnExamen{
% \begin{indicaciones}
%     \item La prueba tiene una duraci\'on  de \textbf{120 minutos.}
%     \item No se permite la utilizaci\'on de dispositivos electr\'onicos (celular, tablets, calculadora, etc.)
%     \item Leer detenidamente los enunciados, \textbf{justificar la respuesta para acceder al puntaje total}.
%     \item La prueba es completamente individual. Los estudiantes deberán dejar sus teléfonos celulares guardados y no podrán tenerlos consigo (por ejemplo, en los bolsillos de su ropa) durante toda la evaluación. En caso de detectarse algún acto de deshonestidad académica o el incumplimiento de esta norma, no podrán seguir rindiendo la prueba y su calificación será \textbf{cero}.
% \end{indicaciones}
% \pagestyle{empty}
% }}

\newcommand{\indicacionesPrueba}{%
\EnExamen{
\begin{indicaciones}
    \item La prueba tiene una duraci\'on  de \textbf{60 minutos.}
    \item Leer detenidamente los enunciados, \textbf{justificar la respuesta} para acceder al puntaje total.
    \item No se permite la utilizaci\'on de dispositivos electr\'onicos (celular, tablets, calculadora, etc.)
    \item La prueba es completamente individual
    \item No se permite portar teléfonos celulares, relojes inteligentes ni otros dispositivos electrónicos durante el examen, aun cuando permanezcan en los bolsillos. Estos dispositivos deberán guardarse en la mochila antes del inicio de la evaluación y esta deberá mantenerse cerrada durante todo el examen.
    \item He leído las indicaciones y acepto que el incumplimiento de estas se penaliza con \textbf{cero}. Firma:...................................

    
\end{indicaciones}
\pagestyle{empty}
}}

\newcommand{\indicacionesExamen}{
\EnExamen{\begin{indicaciones}
    \item La evaluación tiene una duraci\'on  de \textbf{120 minutos.}
    \item No se permite la utilizaci\'on de dispositivos electr\'onicos (celular, tablets, calculadora, etc.)
    %\item Leer detenidamente los enunciados, justificar la respuesta para acceder al puntaje total.
    \item El examen es completamente individual. En caso de detectarse alg\'un acto de deshonestidad acad\'emica no podr\'a seguir rindiendo la prueba y su calificaci\'on ser\'a \textbf{cero}.
\end{indicaciones}
}\pagestyle{empty}
}

%% Imprimir las una hoja A4, dos veces en A5
\newcommand{\imprimirDoble}
{%
\usepackage{pgfmorepages}
\pgfmorepagesloadextralayouts
\pgfpagesuselayout{repeated 2-up}[a4paper,landscape]
\thispagestyle{empty}
}


%Este sí se puede comentar
\newenvironment{solucion}{\color{blue}\proof[Solución]}{\endproof\color{black}}




% \renewcommand{\@autor}{
% ---\\
% \phantom{.} 
% ...}
% \makeatother

%\imprimirDoble
\begin{document}
\encabezado
\datosAlgebra{GR6M}{17/11/2026}
%Poner el grupo y la fecha, o dejar las llaves vacías.
\indicacionesPrueba
\begin{preguntas}


%%%%%%%%%%%%%%%%%%%%%%%%%%%%%%%%%
%%%%%%%%%%%%%%%%%%%%%%%%%%%%%%%%%

\item Sea $\alpha \in \mathbb{R}$. Considere el siguiente sistema de ecuaciones:
    $$
    \left\{
    \begin{alignedat}{4}
        x_1  & - \alpha x_2       & + 4x_3           &= 4 \\
        2x_1 & - (\alpha + 3)x_2  & + (2\alpha + 9)x_3 &= 10 \\
        x_1  & - \alpha x_2       & + (\alpha + 9)x_3  &= 3
    \end{alignedat}
    \right.
    $$
    \begin{enumerate}
        \item Halle el determinante de la matriz de coeficientes, e indique para qué valores de $\alpha$ el sistema tiene solución única. 
        \item Para $\alpha = 3$, halle el conjunto solución. 
    \end{enumerate}
\newpage
%%%%%%%%%%%%%%%%%%%%%%%%%%%%%%%%%%%%%%%
\item Sean los vectores de \(\mathbb{R}^3\):
\[
v_1 = (1, 2, 1), \quad v_2 = (2, 1, -1), \quad v_3 = (1, -4, k),
\]
donde \(k \in \mathbb{R}\) es un parámetro.

\begin{enumerate}
    \item Determine \textbf{para qué valores de \(k\)}, el vector \(v_3\) es combinación lineal de \(v_1\) y \(v_2\).
    \item Determine el \(\operatorname{span}\{v_1, v_2\}\) e interprete geométricamente si el espacio generado es una recta o un plano en \(\mathbb{R}^3\). Justifique.
    \item Dado el vector \(v = (4, 5, 1)\). Determine si el vector \(v\) pertenece a \(\operatorname{span}\{v_1, v_2\}\).
   \item Para el valor de \(k\) hallado en el literal (a), determine si se cumple que \(\operatorname{span}\{v_1, v_2, v_3\} = \operatorname{span}\{v_1, v_2\}\). Justifique su respuesta.
\end{enumerate}
\newpage
%%%%%%%%%%%%%%%%%%%%%%%%%%%%%%%%%%%%%%%%%
\item  Considere S y T subespacios vectoriales $\mathcal{P}_2$ y, $p(x)=ax^2+bx+c$, 
\begin{center}
$S=\{p \in \mathcal{P}_2 : a+b=0\},\qquad
T=\{p\in \mathcal{P}_2 : b+c=0,\ a-c=0\}.$
\end{center}
\begin{itemize}
\item[(a)] Hallar $S\cap T$ de forma explícita.
\item[(b)] Hallar $k\in\R$ tal que $p(x)=(k+1)x^2-3x+(2k-1)\in S\cap T$.
\end{itemize}

\vspace*{30\baselineskip}
%%%%%%%%%%%%%%%%%%%%%%%%%%%%%%%%%%%%%%%
\item (Ejercicio del deber) Analice si el siguiente conjunto es un subespacio vectorial de $\R^{3 \times 3}$:
    \[
    W = \left\{ A \in \R^{3 \times 3} \colon \operatorname{tr}(A) = 0 \right\}
    \] 
\newpage

%%%%%%%%%%%%%%%%%%%%%%%%%%%%%%%%%%%%%%%%%
\item \textbf{Servicios de un taller.} Un taller cobra tarifas fijas por diagnóstico, ajuste y limpieza de un equipo. Un servicio que incluye las tres tareas cuesta \$60. Si el cliente solicita únicamente ajuste y limpieza, paga \$35. Además, el diagnóstico cuesta \$10 más que la limpieza. Determine la tarifa de cada tarea.

\begin{enumerate}
    \item Defina las tres variables, plantee el sistema de ecuaciones lineales que modela la situación y escriba su representación matricial $A\mathbf{x}=b$, identificando $A$, $\mathbf{x}$ y $b$.
    \item Resuelva el sistema e indique la tarifa de las tareas (servicios) detalladas en el problema. 
\end{enumerate}

   

\end{preguntas}

\end{document}